% ===================================================================
%  図表第 10 号 — 海。— 港。
%
%  Ceci n'est PAS une transcription : c'est une traduction, faite en
%  2026, en japonais d'aujourd'hui. Voir texto/ja/10-zuhyo-01.tex
%  pour la consigne, qui vaut pour les seize fichiers.
% ===================================================================

%%K t10-apar-1 apar
\VUsaut{4.0mm}
\VUtitre{15.0pt}{60}{図表第 10 号}
\VUsaut{3.0mm}
\VUpk{11.4pt}{40}{海。--- 港。}
\VUsaut{3.0mm}

%%K t10-01-1 p
1. --- 夏には山へ静けさと涼しさを求めに行く人もあれば、海辺へ行く方を好む人
もある。去年は我々もそうした。\VUgras{大洋} \textsuperscript{(1)} がたいそう好きだ
からである。

%%K t10-02-1 p
2. --- 私はといえば、\VUgras{水平線} \textsuperscript{(2)} まで敷きつめられた
あの果てしない水の面を眺めること、そして巨大な\VUgras{汽船}
\textsuperscript{(3)} がアメリカへ行く旅客を載せて長い\VUgras{航海}に出るのを見る
ことが、この上ない楽しみである。

%%K t10-03-1 p
3. --- だから父は私の好みを知っていて、休みにはいつも我々の大きな
\VUgras{軍港}のそばのたいそう静かな浜を選ぶ。

%%K t10-03-2 p
前に滞在したとき、\VUgras{艦隊}が来てあの巨大な\VUgras{防波堤}
\textsuperscript{(4)} の後ろに錨を下ろした --- 防波堤は\VUgras{軍港の内海}
\textsuperscript{(5)} を閉ざし護っている --- ので、私はいろいろな\VUgras{軍艦}を
心ゆくまで眺めることができた。潜って水中に消える小さな\VUgras{潜水艇}
\textsuperscript{(10)} から、あの巨大な\VUgras{甲鉄艦} \textsuperscript{(9)} まで
である。甲鉄艦は今日まで\VUgras{水雷艇} \textsuperscript{(8)} の襲撃のほかに恐れる
ものがなかった。

%%K t10-tit-2 sub
\VUsaut{3.0mm}
\VUpk{10.2pt}{40}{小舟。}
\VUsaut{3.0mm}

%%K t10-04-1 p
4. --- 水雷艇はとうに厚い甲鉄の弱みをたいそう\VUgras{巧み}に探り当て、そこへ
危うい\VUgras{水雷}を仕掛けることができた。爆ぜれば船を海の底へ送るのである。
それでもやはり潜水艇ほどには恐ろしくない。駆逐艦がたやすく追いつけるから
である。

%%K t10-05-1 p
5. --- 私は速い装甲の\VUgras{巡洋艦} \textsuperscript{(6)} も見た。まことの甲鉄艦
とほとんど同じく刃を通さない。それに\VUgras{兵員輸送船} \textsuperscript{(12)} が
幾隻か、\VUgras{砲艦} \textsuperscript{(7)} が三四隻、最後に\VUgras{護衛艦}
\textsuperscript{(11)} が一隻である。\VUgras{あの艦隊が錨を上げる有様こそ最も見て
面白い。}

%%K t10-06-1 p
6. --- あの鋼鉄の巨きな塊と、商船隊でいまも使う美しい\VUgras{三本檣船}や
瀟洒な\VUgras{帆船} \textsuperscript{(13)} とでは、姿がなんと違うことか。私は
\VUgras{波止場} \textsuperscript{(14)} を散歩しながら、帆をいっぱいに張って港へ
入ってくる帆船をよく見た。もちろん\VUgras{風}が逆であれば帆船はずいぶん損を
する。船渠へ入るには帆をみな畳まねばならず、\VUgras{曳船} \textsuperscript{(16)}
に迎えに行ってもらい、ただの\VUgras{荷足船} \textsuperscript{(17)} のように岸まで
曳いてもらわねばならない。

%%K t10-tit-3 sub
\VUsaut{3.0mm}
\VUpk{10.2pt}{40}{商港。}
\VUsaut{3.0mm}

%%K t10-07-1 p
7. --- 私の言うこの港が面白いのは、巨大な工事で人の手に造った軍港の内海
だけでなく、大きな川の\VUgras{河口}にすばらしい\VUgras{商港}をも備えている
からである。

%%K t10-08-1 p
8. --- 二つの船渠を隔てる地帯には工事が施され、海へ突き出た\VUgras{砲台}の
列と\VUgras{稜堡}とが護っている。\VUgras{城壁} \textsuperscript{(15)} を散歩するには
特別の許しが要る。私は叔父の口ききでたやすくそれを得た。叔父は\VUgras{造船所}
\textsuperscript{(18)} の技師である。それで私は大きな覆屋の下で造られている船を
見に行った。

%%K t10-09-1 p
9. --- 私は甲鉄艦の進水も見た。あの巨きな塊が手を離れるや船台をすべり下り、
川の水に浮かんだとき、群衆じゅうが覚えたあの高ぶりを私は覚えている。

%%K t10-10-1 p
10. --- 造船所へ行くには\VUgras{練兵場} \textsuperscript{(19)} を横切らねばなら
ない --- 守備の兵が毎日そこで調練するのである --- そして鉄の\VUgras{歩み橋}
\textsuperscript{(20)} を渡らねばならない。その橋が練兵場と\VUgras{武器庫}
\textsuperscript{(21)} とをつないでいる。

%%K t10-tit-4 sub
\VUsaut{3.0mm}
\VUpk{10.2pt}{40}{武器庫と船渠。}
\VUsaut{3.0mm}

%%K t10-11-1 p
11. --- 私は叔父とあの大きな施設を見た。中は\VUgras{大砲} \textsuperscript{(22)}、
\VUgras{砲弾} \textsuperscript{(23)}、\VUgras{炸裂弾}でいっぱいである。\VUgras{鍛冶場} \textsuperscript{(24)} も見た。そこで汽船の\VUgras{汽罐}
\textsuperscript{(25)} を造るのである。

%%K t10-12-1 p
12. --- すぐそばが\VUgras{船渠} \textsuperscript{(26)} --- 乾船渠 --- であり、
それにたいそう見るに値する\VUgras{浮船渠} \textsuperscript{(27)} がある。修理中の
船で、私は船の\VUgras{螺旋} \textsuperscript{(28)}、\VUgras{竜骨}
\textsuperscript{(29)}、\VUgras{舵} \textsuperscript{(30)} を間近に見ることができた。

%%K t10-13-1 p
13. --- 商港も軍港に劣らず見るに値し、私は幾時間も波止場で過ごした。そこには
大きな川をさかのぼる\VUgras{外輪船} \textsuperscript{(31)} が繋いである。そして
\VUgras{三脚起重機} \textsuperscript{(32)} が働くのを眺めた。きわめて重いものを
羽根のように持ち上げるのである。

%%K t10-tit-5 sub
\VUsaut{4.0mm}
\VUpk{10.2pt}{40}{水先案内。}
\VUsaut{3.0mm}

%%K t10-14-1 p
14. --- 私はあの\VUgras{遠洋郵船} \textsuperscript{(34)} が\VUgras{旗}
\textsuperscript{(35)} を掲げて泊地を出てゆくのに見とれた。腕のよい
\VUgras{水先案内} \textsuperscript{(36)} が舵を執っていて、\VUgras{浮標}
\textsuperscript{(40)} と\VUgras{標柱}で示された\VUgras{水路}を見事にたどり、
ただ一枚の四角い帆でひどく扱いにくい\VUgras{荷足船} \textsuperscript{(41)} を
よけてゆく。\VUgras{船首} \textsuperscript{(37)} には二等の\VUgras{船室}
\textsuperscript{(39)} が、\VUgras{船尾} \textsuperscript{(38)} には一等の客の広間が
あるのが見分けられた。

%%K t10-15-1 p
15. --- ときには\VUgras{貨物船} \textsuperscript{(42)} の積みこみも見た。
\VUgras{倉庫} \textsuperscript{(43)} と\VUgras{港の停車場} \textsuperscript{(44)} から
運ばれた荷を詰めこむのである。その荷は\VUgras{波止場の鉄道} \textsuperscript{(45)}
の幾列もの汽車が運んできたものである。

%%K t10-tit-6 sub
\VUsaut{3.0mm}
\VUpk{10.2pt}{40}{舟を漕ぐ楽しみ。}
\VUsaut{3.0mm}

%%K t10-16-1 p
16. --- 私はよく\VUgras{内船渠} \textsuperscript{(53)} で舟を漕いだ。小さな
\VUgras{舟} \textsuperscript{(46)} はみなそこに避けているのである。\VUgras{櫂}
\textsuperscript{(47)} を握るのは私の楽しみであった。\VUgras{甲板}
\textsuperscript{(48)} にのぼりもした --- それは\VUgras{帆船} \textsuperscript{(49)} の
甲板である --- \VUgras{索具} \textsuperscript{(52)} を伝って\VUgras{檣}
\textsuperscript{(51)} を登り、\VUgras{帆桁} \textsuperscript{(50)} まで行った。
それからまた\VUgras{小舟} \textsuperscript{(54)} で、重い\VUgras{漁船}
\textsuperscript{(55)} のあいだを漕ぎ続けた。漁船には\VUgras{網} \textsuperscript{(56)}
が干してある。

%%K t10-17-1 p
17. --- \VUgras{波止場} \textsuperscript{(58)} では、ありとあらゆる\VUgras{荷}
\textsuperscript{(59)} が私の前を過ぎていった。\VUgras{木箱} \textsuperscript{(60)} に
入れ、あるいは\VUgras{梱} \textsuperscript{(61)} に束ねてある。それが荷足船の
\VUgras{積荷} \textsuperscript{(63)} をなす。力強い\VUgras{起重機}
\textsuperscript{(62)} がそれを吊り上げ、\VUgras{荷車} \textsuperscript{(64)} に載せる。
そのあいだ\VUgras{運送人}は\VUgras{税関} \textsuperscript{(65)} で申告し税を納めて
いる。

%%K t10-18-1 p
18. --- 最後に\VUgras{回転橋} \textsuperscript{(66)} か\VUgras{閘門}
\textsuperscript{(69)} まで来て、\VUgras{浚渫船} \textsuperscript{(68)} のそばを通り
--- それは\VUgras{運河} \textsuperscript{(67)} を浚っている --- \VUgras{信号柱}
\textsuperscript{(70)} のそばで岸に上がるのであった。

%%K t10-19-1 p
19. --- その堤の反対側に、艦隊の士官を陸へ運ぶ\VUgras{汽艇} \textsuperscript{(71)}
が着く。水兵たちが拍子を合わせて櫂を上げ、舟が\VUgras{波} \textsuperscript{(72)} に
乗って難なく上がり段に着くさまは美しい。ここでこそ\VUgras{潮}が、
\VUgras{浜} \textsuperscript{(73)} での\VUgras{引潮}と\VUgras{満潮}の行き来が、
最もよく見える。

%%K t10-tit-7 sub
\VUsaut{3.0mm}
\VUpk{10.2pt}{40}{士官の集会所。}
\VUsaut{3.0mm}

%%K t10-20-1 p
20. --- 家へ帰るには\VUgras{電車} \textsuperscript{(74)} に乗った。その
\VUgras{停留所} \textsuperscript{(75)} は士官の集会所の前に立つ\VUgras{柱}の
ところである。

%%K t10-20-2 p
集会所の\VUgras{露台} \textsuperscript{(76)} から --- 露台は\VUgras{錨}
\textsuperscript{(77)}、大砲、砲弾で飾ってある --- 私はよく海軍士官の華やかな集りを
見た。剣を佩いた\VUgras{中尉} \textsuperscript{(78)}、\VUgras{軍刀}を佩いた砲兵の
\VUgras{大尉} \textsuperscript{(79)} と\VUgras{歩兵} \textsuperscript{(80)} の大尉で
ある。その後ろには\VUgras{水兵} \textsuperscript{(81)} が控えていて、遠くの様子を
見分けたいときに\VUgras{望遠鏡}を差し出すのである。

%%K t10-21-1 p
21. --- 若い\VUgras{少尉} \textsuperscript{(82)} が\VUgras{艦長} \textsuperscript{(83)}
や\VUgras{提督} \textsuperscript{(84)} から命を受けるのも見た。そのあいだ
\VUgras{兵曹} \textsuperscript{(85)} がそれを艦へ運ぼうと待っていた。

%%K t10-22-1 p
22. --- その露台からの眺めは見事である。浜がすべて見わたせる。浜には
\VUgras{天幕} \textsuperscript{(86)} と\VUgras{着替小屋} \textsuperscript{(87)} が
あり、海水浴の時刻には\VUgras{浴する人} \textsuperscript{(88)} が\VUgras{賭博場}
\textsuperscript{(89)} の前で楽しげに戯れるのが見える。

%%K t10-23-1 p
23. --- 浜のはずれで岸は小さな岩の\VUgras{半島} \textsuperscript{(91)} をなし、
そこへ\VUgras{砕け波} \textsuperscript{(92)} が打ちつける。半島の上には
\VUgras{灯台} \textsuperscript{(93)} が建ち、夜、船人に道を示す。この半島は
たいそう狭い\VUgras{地峡} \textsuperscript{(90)} だけで陸につながっている。

%%K t10-tit-8 sub
\VUsaut{3.0mm}
\VUpk{10.2pt}{40}{海岸の全景。}
\VUsaut{3.0mm}

%%K t10-24-1 p
24. --- \VUgras{断崖} \textsuperscript{(94)} はどこまでも続き、海に刻まれて、
気まぐれな\VUgras{入江} \textsuperscript{(95)} と\VUgras{岬}の連なりを彫り、
それが崖をこの上なく美しく見せている。

%%K t10-24-2 p
\VUgras{台地} \textsuperscript{(96)} は\VUgras{海峡} \textsuperscript{(98)} を
見下ろしていて、その上にたいそう大切な\VUgras{砦} \textsuperscript{(97)} と別荘が
幾つかある。

%%K t10-25-1 p
25. --- その海峡はたいそう危うい\VUgras{島} \textsuperscript{(99)} を大陸から
隔てている。島のまわりは\VUgras{岩礁} \textsuperscript{(100)} と\VUgras{暗礁}で、
小舟はもとより大きな船までがときには難に遭う。救いが早く、\VUgras{救助艇}が
間に合っても、この\VUgras{難破}はやはり恐ろしい。分点のころの\VUgras{大風}では
幾隻もの船が人も荷も共に沈んでいる。

%%K t10-25-2 p
このような隣がありながら、この\VUgras{港}にはやはり船人がよく通う。
\VUsaut{6.0mm}
\VUfilet{22mm}
